\chapter{IgA Nephropathy}
\index{IgA Nephropathy}
\label{sec:IgA Nephropathy}

\section{Overview}
IgA nephropathy (Berger's disease) is the most common primary glomerulonephritis worldwide, with higher prevalence in East Asia and Southern Europe.
\section{Causes}
This condition happens because of deposition of galactose-deficient IgA1 in the mesangium, triggering an immune response, mesangial growth, and glomerular injury.     
\section{Risk Factors}

\begin{itemize}[noitemsep]
  \item Genetic predisposition and family history
  \item Advancing age
  \item Lifestyle factors (diet, exercise, smoking, alcohol)
  \item Environmental exposures
  \item Underlying medical conditions
  \item Medication use
\end{itemize}
\section{Signs and Symptoms}

\subsection{Common symptoms}
\begin{itemize}[noitemsep]
  \item You may experience episodic gross hematuria (often concurrent with upper respiratory infections---"synpharyngitic hematuria"), microscopic hematuria, or proteinuria
  \item Some present with nephrotic syndrome or rapidly progressive glomerulonephritis.
  \item Pain or discomfort in the affected area
  \end{itemize}

\subsection{Emergency warning signs}
\begin{itemize}[noitemsep]
  \item Severe or worsening symptoms of iga nephropathy require immediate medical evaluation
\end{itemize}
\section{Progression and Stages}
The condition may progress through different stages of severity over time. Early stages often involve mild or subtle symptoms that may be overlooked. Moderate stages involve more noticeable symptoms affecting daily function. Advanced stages may involve significant impairment and complications.
\section{Complications}
\begin{itemize}[noitemsep]
  \item Disease progression leading to worsening symptoms
  \item Secondary infections or injuries
  \item Side effects of treatment
  \item Reduced quality of life
  \item Emotional and psychological impact
\end{itemize}
\section{Diagnosis}
\begin{itemize}[noitemsep]
  \item To make the diagnosis, doctors look for kidney biopsy showing mesangial IgA deposition on immunofluorescence
  \item The Oxford Classification (MEST-C scoring) guides outlook.
  \item Comprehensive medical history and physical examination
  \item Laboratory tests to assess organ function
  \item Imaging studies to visualize affected structures
  \item Specialized testing depending on the affected system
  \item Consultation with relevant specialists
\end{itemize}
\section{Prevention}
\begin{itemize}[noitemsep]
  \item Maintain a healthy lifestyle with balanced diet and exercise
  \item Avoid known risk factors
  \item Attend regular medical check-ups for early detection
  \item Follow recommended screening guidelines
\end{itemize}
\section{Treatment Options}
Neflgentan (sodium-glucose cotransporter 2 inhibitor) is an emerging therapy. Medications to manage symptoms and address underlying mechanisms Lifestyle modifications and self-care measures Physical therapy and rehabilitation Surgical intervention when indicated Regular monitoring and follow-up care Patient education and support services

\begin{itemize}[noitemsep]
  \item Medications to manage symptoms and address underlying mechanisms
  \item Lifestyle modifications and self-care measures
  \item Physical therapy and rehabilitation
  \item Surgical intervention when indicated
  \item Regular monitoring and follow-up care
\end{itemize}
\section{Living with It}
\begin{itemize}[noitemsep]
  \item The right treatment depends on severity: RAAS blockade for proteinuria, SGLT2 inhibitors, and immunosuppression (corticosteroids) for high-risk patients
  \item Supportive care includes BP control and smoking cessation.
  \item Follow your treatment plan as prescribed
  \item Maintain regular follow-up with your healthcare provider
  \item Make necessary lifestyle adjustments
  \item Seek support from family, friends, or support groups
  \item Stay informed about your condition
\end{itemize}
\section{Outlook}
\begin{itemize}[noitemsep]
  \item The outlook varies from person to person
  \item Up to 30\% progress to ESRD within 20 years.
\end{itemize}
